% 矮行星（综述）
% license CCBYSA3
% type Wiki

本文根据 CC-BY-SA 协议转载翻译自维基百科\href{https://en.wikipedia.org/wiki/Dwarf_planet}{相关文章}。

\begin{figure}[ht]
\centering
\includegraphics[width=6cm]{./figures/d21611a5210b90a2.png}
\caption{} \label{fig_Dwarf_1}
\end{figure}
“矮行星是一类直接绕太阳运行、具有行星质量但体积较小的天体。它们的质量足以使自身达到由重力塑形的近球形，但不足以像太阳系中八大经典行星那样在轨道上实现支配地位。矮行星的原型是冥王星，它曾被视为一颗行星长达数十年，直到 2006 年‘矮行星’概念被正式采用。许多行星地质学家认为矮行星与具有行星质量的卫星都应被视为行星\(^\text{[1]}\)，但自 2006 年以来，国际天文学联合会（IAU）和许多天文学家已将它们排除在行星名录之外。

矮行星可能具有地质活动能力——这一预期在 2015 年“黎明号”对谷神星的探测和“新视野号”对冥王星的探测中得到证实。因此，它们成为行星地质学家特别关注的对象。

天文学家普遍同意至少有九个最大的候选体属于矮行星——按直径从大到小大致排列为：冥王星、厄里斯、妊神星、鸟神星、公公星、夸欧尔、塞德娜、谷神星和奥库斯。对于第十大候选体萨拉西娅仍存在相当的不确定性，因此可将其视为临界案例。在这十个天体中，已有两个被航天器访问过（冥王星和谷神星），另有七个至少拥有一颗已知卫星（厄里斯、妊神星、鸟神星、公公星、夸欧尔、奥库斯和萨拉西娅），这使得它们的质量得以确定，从而可估算其密度。质量和密度又可用于构建地球物理模型，以尝试判断这些世界的性质。只有塞德娜既未被探测器访问，也没有已知卫星，因此其质量难以准确估计。一些天文学家还将许多更小的天体也纳入矮行星的范畴\(^\text{[2]}\)，但目前并无共识认为这些小天体很可能是矮行星。”
\subsection{概念的历史}
\begin{figure}[ht]
\centering
\includegraphics[width=6cm]{./figures/0ce0512cdfe2e0c4.png}
\caption{冥王星及其卫星卡戎的近真彩色影像。二者的间距按比例绘制} \label{fig_Dwarf_2}
\end{figure}
\begin{figure}[ht]
\centering
\includegraphics[width=6cm]{./figures/6ae4d90716be2481.png}
\caption{4 号灶神星（Vesta），一度曾是矮行星的天体\(^\text{[3]}\)} \label{fig_Dwarf_3}
\end{figure}
自1801年起，天文学家发现了谷神星以及位于火星与木星之间的其他天体，这些天体在随后的数十年间都被视为行星。自那时直到大约1851年，当行星数量已达到23颗时，天文学家开始使用“asteroid”（来自希腊语，意为“类星的”或“星形的”）一词来指代较小的天体，并开始将它们区分为小行星而非大行星。\(^\text{[4]}\)

随着1930年冥王星的发现，大多数天文学家认为太阳系包含九颗主要行星，以及成千上万的显著更小的天体（小行星与彗星）。在将近50年的时间里，人们曾认为冥王星比水星更大，\(^\text{[5][6]}\)但在1978年冥王星的卫星卡戎被发现后，天文学家得以精确测量冥王星的质量，并确定其远小于最初的估计。\(^\text{[7]}\)其质量大约只有水星的二十分之一，使冥王星成为迄今最小的行星。尽管冥王星仍比小行星带中最大的天体谷神星大十倍以上，但其质量却只有地球月球的五分之一。\(^\text{[8]}\)此外，由于其具有一些异常特征，如较大的轨道离心率与较高的轨道倾角，愈发明显冥王星与其他行星类型不同。\(^\text{[9]}\)

在1990年代，天文学家开始在与冥王星相同的空间区域（即如今所称的柯伊伯带）以及更远处发现更多天体。\(^\text{[10]}\)其中许多与冥王星共享关键的轨道特征，冥王星开始被视为一类新天体的最大成员，即冥族（plutinos）。于是变得明显：要么这些较大的天体也需要被归类为行星，要么冥王星需要像谷神星在更多小行星被发现后那样被重新分类。\(^\text{[11]}\) 这使得一些天文学家不再将冥王星称为行星。多个术语包括“subplanet”（次行星）与“planetoid”（类行星体）等，被用于指代如今所谓的矮行星。\(^\text{[12][13]}\)天文学家也确信会有更多与冥王星大小相当的天体被发现，如果冥王星继续被视为行星，行星数量将会快速增长。\(^\text{[14]}\)

海王星外天体厄里斯（当时称为 2003 UB313）于 2005 年 1 月被发现；\(^\text{[15]}\)当时认为其尺寸略大于冥王星，因此一些报道非正式地称其为“第十颗行星”。\(^\text{[16]}\)因此，该问题在 2006 年 8 月的国际天文学联合会（IAU）大会上成为激烈争论的焦点。\(^\text{[17]}\)IAU 的最初草案提议将卡戎、厄里斯与谷神星列入行星名单。由于许多天文学家反对这一方案，乌拉圭天文学家 Julio Ángel Fernández 与 Gonzalo Tancredi 制定了另一项替代方案：他们提出设立一个中间类别，用于归类那些质量足以达到流体静力平衡（呈球形），但并未清除其轨道中微行星的天体。除了将卡戎从名单中移除之外，该新方案也将冥王星、谷神星与厄里斯排除在行星名单之外，因为它们均未清除其所在轨道。\(^\text{[18]}\)

尽管有人对此分类在系外行星中的适用性提出了担忧，\(^\text{[19]}\)该问题仍未得到解决；提议是等到能够观测到与矮行星体积相当的系外天体时再作决定。\(^\text{[18]}\)

在 IAU 给出“矮行星”定义的直接后续中，一些科学家表达了他们对 IAU 决议的不满。\(^\text{[20]}\)抗议活动包括汽车保险杠贴纸与 T 恤标语。\(^\text{[21]}\)厄里斯的发现者 Mike Brown 则赞同将行星数量减少至八颗的做法。\(^\text{[22]}\)

NASA 在 2006 年宣布将采用 IAU 制定的新指南。\(^\text{[23]}\)NASA 冥王星任务的负责人 Alan Stern 不认可当前 IAU 的行星定义，无论是将矮行星排除在行星类别之外，还是依据轨道特征（而非天体的固有特征）来定义矮行星。\(^\text{[24]}\)因此，他在 2011 年仍将冥王星称为行星，\(^\text{[25]}\)并将其他可能的矮行星，如谷神星与厄里斯，以及更大的卫星，也视作额外的行星。\(^\text{[26]}\)在 IAU 发布定义的数年前，他便使用轨道特征将“überplanets”（主导性的八颗行星）与“unterplanets”（矮行星）区分开来，并认为二者皆为“行星”。\(^\text{[27]}\)
\subsection{名称}
\begin{figure}[ht]
\centering
\includegraphics[width=6cm]{./figures/46cf07ad5d90f072.png}
\caption{欧拉图：展示国际天文学联合会（IAU）执行委员会对太阳系（不含太阳）内天体类别的构想} \label{fig_Dwarf_4}
\end{figure}
对大型亚行星天体的称呼包括 dwarf planet（矮行星）、planetoid（类行星，更泛化的术语）、meso-planet（中型行星，狭义上指介于水星与谷神星之间尺寸的天体）、quasi-planet（准行星），以及（在海王星外区域）plutoid（冥族天体）。然而，dwarf planet 这一术语最初是用于指代“最小的行星”，而不是“最大的亚行星”，并且至今仍以这种含义被许多行星地质学家使用。

dwarf planet 作为 asteroid（小行星）的同义词可以追溯至少到 1838 年。它也曾作为 giant planet（巨行星）的反义词，被用于包括类地行星——水星、金星、地球和火星——以及月球在内的天体。\(^\text{[28]}\)在 2006 年的分类中，IAU 规定矮行星不应被视为行星。对于 IAU 所界定的最大亚行星天体，其他不带有类似歧义或用法冲突的术语包括 quasi-planet\(^\text{[29]}\)与较早的 planetoid（“具有行星外形之物”）。\(^\text{[30]}\)Michael E.~Brown 指出，planetoid “是一个完全没问题的词汇，使用多年”，而用 dwarf planet 去指代“非行星”则是“很傻”（dumb），但这是因为 IAU 第三分部大会试图通过第二项决议让冥王星重新成为行星。\(^\text{[31]}\)事实上，5A 号决议草案原本将这些中等大小的天体称为 planetoids，\(^\text{[32][33]}\)但大会一致投票将其名称改为 dwarf planet。\(^\text{[34]}\)第二项决议 5B 则将矮行星定义为行星的一种子类型（与 Stern 最初的设想一致），与其它八颗行星区分为“classical planets”（经典行星）。在这种安排下，被否决的提案所列的 12 颗“行星”本可以通过“八颗经典行星 + 四颗矮行星”这一区分得以保留。然而，决议 5B 在通过 5A 的同一场会议中被否决了。\(^\text{[31]}\)由于决议 5B 未获通过，从语义上造成了“矮行星并非行星”这一不一致，因此出现了 nanoplanet、subplanet 等替代名词的讨论，但 IAU 的 CSBN 并未就更名达成共识。\(^\text{[35]}\)

在大多数语言中，“矮行星”都被更或少直译为对应词语：法语 planète naine、西班牙语 planeta enano、德语 Zwergplanet、俄语 karlikovaya planeta（карликовая планета）、阿拉伯语 kaukab qazm（كوكب قزم）、中文 矮行星、韩语 왜소행성／waesohangseong（矮小行星）或 왜행성／waehangseong（矮行星）。但在日语中，它们被称为 准惑星（junwakusei），意为“准行星”或“近似行星”（peneplanet，pene- 表示“几乎”）。

2006 年的 IAU 第 6a 号决议\(^\text{[36]}\)将冥王星认定为“新一类海王星外天体的原型”。该类别的名称与精确定性当时尚未指定，而是留待 IAU 日后确定；在该决议出台前的讨论中，这一类别的成员曾被称为 plutons 与 plutonian objects，但这两个名称均未被采纳，可能是因为地质学家反对这种命名会与地质学中的 pluton（深成岩体）混淆。\(^\text{[34]}\)

2008 年 6 月 11 日，IAU 执行委员会宣布了一个新术语 plutoid，并给出了定义：所有海王星外的矮行星皆为 plutoid。\(^\text{[37]}\)然而，IAU 的其他部门拒绝了这一术语：

……部分由于电子邮件沟通失误，WG-PSN（行星系统命名工作组）并未参与“plutoid”一词的选择。……事实上，该工作组在执行委员会会议之后所进行的投票否决了这一特定术语的使用……\(^\text{[35]}\)

“plutoid”这一类别继承了先前根据地质性质在“类地矮行星”谷神星与外太阳系“冰矮星”之间所做的区分。\(^\text{[38]}\)这一区分是三分太阳系概念的一部分：即太阳系由内侧类地行星、中部巨行星、以及外侧冰矮星构成，而冥王星是该外侧组的代表成员。\(^\text{[39]}\)“ice dwarf（冰矮星）”过去也曾用作海王星外所有小行星（TNOs）的总称，或指外太阳系的“冰质小行星”；其中一种尝试性的定义是：冰矮星“比典型彗核更大，且比典型小行星更富含冰”。\(^\text{[40]}\)

自 Dawn 号任务以来，人们已经认识到谷神星是一个地质上富冰的天体，并可能起源于外太阳系。\(^\text{[41][42]}\)因此谷神星也被称作冰矮星。\(^\text{[43]}\)\subsection{分类标准}
\begin{figure}[ht]
\centering
\includegraphics[width=10cm]{./figures/cf2087e9444efd75.png}
\caption{} \label{fig_Dwarf_5}
\end{figure}
“矮行星”这一类别的出现，源于关于“何谓行星”这一概念究竟应依据动力学标准还是地球物理学标准而产生的冲突。就太阳系动力学而言，最重要的区分在于：一类天体能够在其轨道邻域内表现出引力支配地位（从水星到海王星），而另一类天体则不能（例如小行星与柯伊伯带天体）。一个天体在其质量达到某一程度时，其地幔会在自身重量作用下变得具可塑性，从而使天体获得圆形外貌，此时它可能具备动力学意义上的“行星形”地质活动。由于获得圆形所需的质量远低于实现轨道邻域引力清空所需的质量，因此存在这样一类天体：其质量足以具备“类世界的外貌”以及类似行星的地质活动，但不足以清除其轨道邻域。例如小行星带中的谷神星，以及柯伊伯带中的冥王星。\(^\text{[47]}\)

动力学研究者通常倾向将“引力支配轨道邻域”作为行星的判定标准，因为从他们的角度看，将质量较小的天体与其邻域群体归为一类更为合理，例如将谷神星视为一颗较大的小行星，将冥王星视为一颗较大的柯伊伯带天体。\(^\text{[48][49]}\)地球科学家则通常以“是否达到流体静力平衡形状”作为判定阈值，因为从他们的视角来看，像谷神星这样拥有内部驱动地质活动的天体，与拥有行星地质的经典行星（如火星）更为相似，而与缺乏内部地质驱动的小型小行星并不相同。正因为存在这一中间层级，才有必要设立“矮行星”这一类别以描述该类型天体。\(^\text{[47]}\)
\subsubsection{轨道支配性}
Alan Stern 与 Harold F. Levison 在 2000 年提出了参数 Λ（大写 lambda），用以表示一次遭遇导致轨道发生给定偏转的概率。\(^\text{[27]}\)在 Stern 的模型中，该参数的取值与天体质量的平方成正比，并与轨道周期成反比。该值可用于估计天体清除其轨道邻域的能力，其中  Λ > 1 表示该天体最终能够清除其轨道区域。研究发现，从最小的类地行星到最大的一批小行星与柯伊伯带天体之间，Λ 的数值存在五个数量级的差距。\(^\text{[44]}\)

利用这一参数，Steven Soter 以及其他天文学家主张依据天体“无法清除其轨道邻域”来区分行星与矮行星：行星能够通过碰撞、俘获或引力扰动清除其轨道附近的小型天体（或建立轨道共振以防止碰撞）；而矮行星由于质量不足，无法完成这一过程。\(^\text{[27]}\)Soter 进一步提出了一个称为“行星判别式”的参数，记作 µ（mu），其代表轨道区域实际清洁度的实验度量（µ 的计算方式为：候选天体的质量除以与其共享轨道区的所有其他天体的总质量），其中  µ > 100 被视为已清除轨道邻域。\(^\text{[44]}\)

Jean-Luc Margot 对 Stern 与 Levison 的概念进行了改进，提出了一个类似的参数 Π（大写 Pi）。\(^\text{[46]}\)该参数基于理论推导，避免了 Λ 中所依赖的经验性数据。Π > 1 即被视为行星，同样地，行星与矮行星之间在 Π 的取值上也存在若干个数量级的差距。

此外，还有其他一些方案试图区分行星与矮行星，\(^\text{[20]}\)但 2006 年 IAU 的定义采用的正是这一概念。\(^\text{[34]}\)
\subsubsection{流体静力平衡}
当天体的自引力产生足够的内部压力时，其物质会进入塑性状态，而足够的塑性会使高起伏下沉、低洼区域被填平，这一过程称为重力松弛（gravitational relaxation）。**
直径小于数公里的天体主要受非重力因素主导，往往呈不规则形状，并可能是碎石堆（rubble piles）。更大的天体中，引力开始变得重要但尚未完全占优势，此时它们呈现出“马铃薯状”。随着天体质量的增加，其内部压力进一步升高，结构更为坚固，形状也随之更加圆滑，直到内部压力足以压过其抗压强度，使天体进入流体静力平衡（hydrostatic equilibrium）。此时，天体在其自转和潮汐效应允许的条件下，已圆整到极限，呈现椭球体形状。这就是“矮行星”的判定界限。\(^\text{[50]}\)

如果一个天体处于流体静力平衡，那么其表层若覆盖一层全球性的液体，该液体的表面形状（除去陨坑裂隙等小尺度地形）将与该天体整体的形状一致。不自转的天体将呈球形，而自转天体呈椭球形。自转越快，天体越扁甚至可能呈不等轴形。若此类自转天体被加热直至熔融，其形状也不会改变。一个极端的例子是 Haumea，它因快速自转可能呈不等轴形，其最长轴长度约为极半径的两倍。

如果天体拥有一个质量巨大的近伴，则潮汐作用会逐渐减慢其自转，直至其达到潮汐锁定（tidally locked）；即始终以同一面朝向其伴星。潮汐锁定的天体也呈不等轴形，但有时仅有轻微差异。地球的月球处于潮汐锁定状态，所有气体巨行星的圆形卫星亦如此。冥王星与其卫星卡戎彼此潮汐锁定；阋神星（Eris）与其卫星迪斯诺米亚（Dysnomia）亦如此；Orcus 与 Vanth 很可能也处于相同状态。\(^\text{[citation needed]}\)

矮行星并无特定的大小或质量上限，因为这些并非其定义依据。实际上也不存在明确的上限：例如，在太阳系的极外区域，一个质量超过水星但尚未清除其轨道区的天体，也将被归类为矮行星而非行星。事实上，Mike Brown 曾试图寻找这样的天体。\(^\text{[51]}\)

下限由天体达到并保持流体静力平衡的要求决定，但具体大小或质量阈值取决于其成分与热演化历史，而非简单的质量。IAU 2006 年新闻稿\(^\text{[52]}\)的问答部分估计，质量超过 (0.5\times 10^{21},\mathrm{kg}) 且半径大于 (400,\mathrm{km}) 的天体“通常”会处于流体静力平衡状态（形状……通常由自引力决定），但所有边缘个案都必须通过观测来确定。\(^\text{[52]}\)

这一估计与截至 2019 年对海王星外完整固体天体的判断相近，其中 Salacia（(r = 423\pm 11,\mathrm{km})，(m = (0.492\pm 0.007)\times 10^{21},\mathrm{kg})）在 2006 年问答标准与后续研究中均被视为边界案例；而 Orcus 稍高于该预期上限。\(^\text{[53]}\)就已测质量的天体而言，没有其他对象接近上述质量上限，但一些未测质量的天体接近所预期的尺寸下限。\(^\text{[citation needed]}\)
\subsection{矮行星的总体分布}
\begin{figure}[ht]
\centering
\includegraphics[width=8cm]{./figures/c883c7e09b0f9dea.png}
\caption{对比多种直径大于 700,km 的外海王星天体（trans-Neptunian objects）的尺寸、反照率与颜色。深色弧段表示这些天体尺寸的不确定范围。} \label{fig_Dwarf_6}
\end{figure}
“虽然矮行星的定义是明确的，但关于某个特定的海王星外天体（TNO）是否足够巨大且具有足够的延展性，使其能够被自身的引力场塑形，其证据常常并不确凿。在某些情形下，对 IAU 判据的解释也仍存在悬而未决的问题。因此，目前符合流体静力平衡判据的已确认海王星外天体的数量仍不确定。

在 2006 年 IAU 通过矮行星类别之前的讨论中所重点考虑的三个天体——谷神星（Ceres）、冥王星（Pluto）与阋神星（Eris）——通常都被视为矮行星，包括那些仍将矮行星归类为行星的天文学家。它们之中只有冥王星被充分观测，从而能够验证其当前形状符合流体静力平衡的预期,\(^\text{[54]}\)。谷神星接近平衡，但仍存在某些尚未解释的引力异常,\(^\text{[55]}\)。阋神星通常被视为矮行星，因为它的质量大于冥王星。”

“按发现顺序，这三个天体分别为：

\begin{enumerate}
\item \textbf{谷神星（Ceres）}——1801 年 1 月 1 日发现，1 月 24 日公布，比海王星早 45 年。被视为行星长达半个世纪，之后被重新分类为小行星。自 2006 年 8 月 24 日通过第 5A 号决议以来，被 IAU 视为矮行星。
\item \textbf{冥王星（Pluto）}——1930 年 2 月 18 日发现，3 月 13 日公布。作为行星被认可了 76 年。于 2006 年 8 月 24 日被 IAU 通过第 6A 号决议明确重新分类为矮行星,\(^\text{[56]}\)。已知五颗卫星。
\item \textbf{阋神星（Eris，2003,UB_{313}）}——2005 年 1 月 5 日发现，7 月 29 日公布。媒体报道中被称为“第十颗行星”。自 2006 年 8 月 24 日 IAU 通过第 5A 号决议以来被视为矮行星，并于同年 9 月 13 日由 IAU 矮行星命名委员会正式命名。已知一颗卫星。\\\\
IAU 仅制定了哪些委员会负责命名可能的矮行星的指导方针：任何绝对星等亮于 +1 的、尚未命名的海王星外天体（因此在最大几何反照率为 1 的极端情况下，其最小直径为 838,km）,\(^\text{[57]}\)，将由由小行星中心与 IAU 行星工作组联合组成的委员会负责命名,[37]。当时（并且截至 2023 年仍然如此），唯一达到这一阈值的天体是 Haumea 与 Makemake。它们一般被认为是矮行星，尽管尚未证明它们处于流体静力平衡状态，而且对于 Haumea 的判断存在一些分歧,\(^\text{[58][59]}\)。
\item \textbf{Haumea（2003,EL_{61}）}——由 Brown 等人于 2004 年 12 月 28 日发现，并由 Ortiz 等人于 2005 年 7 月 27 日公布。2008 年 9 月 17 日由 IAU 矮行星命名委员会命名。已知两颗卫星与一条行星环。
\item \textbf{Makemake（2005,FY_{9}）}——2005 年 3 月 31 日发现，并于 7 月 29 日公布。2008 年 7 月 11 日由 IAU 矮行星命名委员会命名。已知一颗卫星。\\\\
上述五个天体——2006 年讨论的三个（冥王星、谷神星与阋神星）加上 2008 年命名的两个（Haumea 与 Makemake）——通常被呈现为太阳系的矮行星，尽管限制因素（反照率）并不是定义矮行星的标准,\(^\text{[60]}\)。\\\\
天文学界通常也将其他更大的海王星外天体视为矮行星,\(^\text{[61]}\)。至少还有四个天体满足 Brown、Tancredi 等人、Grundy 等人以及 Emery 等人提出的矮行星判据的初步条件，并且也被天文学家普遍称为矮行星：”
\item \textbf{Quaoar（2002,LM_{60}）}——于 2002 年 6 月 5 日发现，并于当年 10 月 7 日公布。已知一颗卫星与两条行星环。
\item \textbf{Sedna（2003,VB_{12}）}——于 2003 年 11 月 14 日发现，并于 2004 年 3 月 15 日公布。
\item \textbf{Orcus（2004,DW）}——于 2004 年 2 月 17 日发现，并于两天后公布。已知一颗卫星。
\item \textbf{Gonggong（2007,OR_{10}）}——于 2007 年 7 月 17 日发现，并于 2009 年 1 月公布。已知一颗卫星。\\\\
例如，JPL/NASA 在 2016 年的观测之后将 Gonggong 称为矮行星,\(^\text{[62]}\)；而西南研究院的 Simon Porter 在 2018 年提到“八大海王星外矮行星”，指的是冥王星、阋神星、Haumea、Makemake、Gonggong、Quaoar、Sedna 与 Orcus,\(^\text{[63]}\)。IAU 本身也在 2022–2023 年度报告中将 Quaoar 称为矮行星,\(^\text{[64]}\)。\\\\
更多天体也被提出作为候选，例如 Brown 提出的 Salacia 与 Máni，Tancredi 等人提出的 Varuna 与 Ixion，以及 Sheppard 等人提出的 Chiminigagua,\(^\text{[65]}\)。多数体积较大的天体拥有卫星，这使得能够测定其质量并进一步得到密度，从而帮助判断其是否可能为矮行星。目前已知最大但无卫星的海王星外天体包括 Sedna、Máni、Aya 和 Ixion。尤其是 Salacia，其质量与直径已知，使其在 IAU 2006 年的问答指引中被视为临界个案,\(^\text{[citation needed]}\)。”
\item \textbf{Salacia（2004,SB_{60}）}——于 2004 年 9 月 22 日发现。已知一颗卫星。
\end{enumerate}
在命名 Makemake 和 Haumea 的时期，人们认为具有冰质核心的海王星外天体（TNO）只需要约 400 km 的直径（约为地球大小的 3\%）——大致相当于卫星米马斯（最小的呈球形的卫星）和普罗透斯（最大但并未达到球形的卫星）的尺寸——即可松弛至引力平衡状态。,\(^\text{[66]}\)研究人员认为，在柯伊伯带中可能存在约 200 个这样的天体，而在其更远处可能还有数千个。,\(^\text{[66][67][68]}\)这也是冥王星最初被重新分类的重要原因之一（即保持“行星”的数量在可接受范围内）。随后研究对这一设想提出了质疑：在柯伊伯带及其更远处的典型环境中，如此小的天体是否能够达到或维持引力平衡仍存疑。,\(^\text{[citation needed]}\)

一些天文学家陆续将若干天体认定为矮行星或可能成为矮行星的候选者。2008 年，Tancredi 等人建议 IAU 正式将 Orcus、Sedna 和 Quaoar 接纳为矮行星（当时 Gonggong 尚未被发现），但 IAU 当时并未处理此事，之后也未进一步讨论。,\(^\text{[69]}\) Tancredi 还认为另外五个海王星外天体 Varuna、Ixion、Achlys、Goibniu 和 Aya 也极有可能是矮行星。自 2011 年起，Brown 基于估计尺寸维护了一份包含数百个候选天体的列表，范围从“几乎确定”为矮行星到“可能”为矮行星。,\(^\text{[70]}\)截至 2019 年 9 月 13 日，Brown 的列表中已有十个海王星外天体的直径被认为大于 900 km（包括 IAU 已命名的四个矮行星加上 Gonggong、Quaoar、Sedna、Orcus、Máni 和 Salacia），被视为“几乎确定”为矮行星；另有 16 个，直径大于 600 km，被视为“高度可能”。,\(^\text{[67]}\)值得注意的是，Gonggong 的直径（1230±50 km）可能比冥王星的球形卫星卡戎（1212 km）还要更大。

但在 2019 年，Grundy 等人基于对 Gǃkúnǁʼhòmdímà 的研究提出，直径小于约 900–1000 km 的暗色、低密度天体（例如 Salacia 和 Varda）从未完全坍塌成致密的行星状固体，而是保留了形成时期的内部孔隙结构（若如此，它们便不能成为矮行星）。他们接受反照率更高（albedo > ≈0.2）,\(^\text{[71]}\)或密度更大（> ≈1.4 g/cc） 的 Orcus 和 Quaoar 很可能已经成为完全固体：,\(^\text{[53]}\)

> Orcus 和 Charon 很可能曾经熔融并发生分异，其较高的密度以及光谱显示其表面由相对纯净的 H₂O 冰组成。而 Gǃkúnǁʼhòmdímà、55637（Uni）、Varda 和 Salacia 的较低反照率与密度则暗示它们从未分异，或即使发生过分异，也仅限于它们的深部内部，并非涉及表层的完全熔融与翻覆。即便内部变暖并塌陷，它们的表面仍可能保持非常寒冷且未被压实。内部挥发物的释放可进一步帮助向外输送热量，从而限制其内部坍缩的程度。一个具有寒冷、相对原始表面但内部部分塌缩的天体，应该呈现非常独特的表面地质特征，例如大量逆冲断层，以反映其内部压缩收缩导致的表面积减小。,\(^\text{[53]}\)

之后发现 Salacia 的密度略高，在不确定度范围内可与 Orcus 相比，尽管其表面依然非常暗。尽管如此，Grundy 等人称其为“矮行星大小”，而将 Orcus 直接称为矮行星。,\(^\text{[72]}\)对 Varda 的后续研究显示其密度也可能较高，但较低密度的可能性仍无法排除。,\(^\text{[73]}\)

2023 年，Emery 等人撰文指出，詹姆斯·韦伯空间望远镜（JWST）在 2022 年进行的近红外光谱观测显示，Sedna、Gonggong 和 Quaoar 可能经历了内部熔融、分异和化学演化，与较大的矮行星 Pluto、Eris、Haumea 和 Makemake 相似，但不同于“所有更小的柯伊伯带天体”。原因在于它们表面存在轻烃（如乙烷、乙炔和乙烯），这暗示甲烷在持续补给，而甲烷最可能来源于内部地球化学过程。另一方面，Sedna、Gonggong 和 Quaoar 表面的 CO 和 CO₂ 含量都很低，与 Pluto、Eris 和 Makemake 相似，但与较小天体形成对比。这表明，在海王星外区域，矮行星的尺度下限应为直径约 ~900 km（因此仅包括 Pluto、Eris、Haumea、Makemake、Gonggong、Quaoar、Orcus 和 Sedna），甚至 Salacia 也可能并非矮行星。,\(^\text{[74]}\)

2023 年对 Máni 的研究显示，它可能具有一个极大的撞击坑，其深度占其直径的 5.7％，比例上大于 Vesta 的 Rheasilvia 撞击盆地——这也是 Vesta 今天通常不被视为矮行星的原因。,\(^\text{[75]}\)

2024 年，Kiss 等人发现 Quaoar 具有与其当前自转速度不相容的椭球形状。他们假设 Quaoar 早期自转很快，处于流体静力平衡，但在其因卫星 Weywot 的潮汐作用而自转减慢时，其形状被“冻结”下来并未改变。,\(^\text{[76]}\)若如此，则类似于土星卫星 Iapetus 的情况：其形状对于当前自转而言过于扁平。,\(^\text{[77][78]}\)Iapetus 通常仍被视为行星质量的卫星，尽管并非总是如此。,\(^\text{[47][79]}\)
\subsubsection{最可能的矮行星}
\begin{figure}[ht]
\centering
\includegraphics[width=10cm]{./figures/efa2607b23abe92e.png}
\caption{} \label{fig_Dwarf_7}
\end{figure}
除 Salacia 外，下表中的海王星外天体（TNO）均被 Brown、Tancredi 等、Grundy 等以及 Emery 等学者一致认为是可能的矮行星，或非常接近矮行星。之所以将 Salacia 列入，是因为它是目前在总体上尚未达成矮行星共识的最大 TNO；按多项标准来看它都属于临界天体，因此在表中以斜体标示。作为比较，Pluto 的卫星 Charon 也被列入——它在 2006 年被 IAU 提议为矮行星。在这些天体中，凡绝对星等亮于 +1 的，就达到了 IAU 行星–小行星联合命名委员会设定的命名阈值，这些对象（以及 IAU 自 2006 年首次讨论该概念以来就被视为矮行星的 Ceres）都在表中被特别标示。

所列矮行星的质量均为其系统质量（若其拥有卫星），但 Pluto 与 Orcus 为例外。
\begin{figure}[ht]
\centering
\includegraphics[width=14.25cm]{./figures/d4438e91a796a3de.png}
\caption{} \label{fig_Dwarf_9}
\end{figure}
\begin{figure}[ht]
\centering
\includegraphics[width=14.25cm]{./figures/ea3e176661d073a4.png}
\caption{} \label{fig_Dwarf_8}
\end{figure}
\subsubsection{符号}
Ceres（⚳）和 Pluto（♇）在被发现时被视为行星，因此获得了行星符号。等到其他天体被发现时，行星符号在天文学界已经基本不再使用。Unicode 目前包含 Quaoar（🝾）、Sedna（⯲）、Orcus（🝿）、Haumea（🝻）、Eris（⯰）、Makemake（🝼）和 Gonggong（🝽）的符号，这些符号主要由占星家使用；它们均由马萨诸塞州的软件工程师 Denis Moskowitz 设计。[82][83][84]NASA 在将 Pluto 作为矮行星提及时，曾使用过 Moskowitz 设计的 Haumea、Eris、Makemake 的符号，以及传统的冥王星占星符号 ⯓。[83][85]较小天体的符号则尚未普遍确立；Unicode 的一份提案中记录了Moskowitz 为 Salacia 设计的符号 。[86]Moskowitz 还为 Charon 设计了另一个符号 ⯕。[87]
\subsection{探测}
\begin{figure}[ht]
\centering
\includegraphics[width=8cm]{./figures/f9e538f327c039ff.png}
\caption{由美国国家航空航天局“黎明号”探测器拍摄的矮行星谷神星} \label{fig_Dwarf_10}
\end{figure}
截至2025年，只有两项任务曾近距离针对并探测过矮行星。2015年3月6日，“黎明号”（Dawn）探测器进入谷神星轨道，成为首个访问矮行星的航天器[88]。2015年7月14日，“新视野号”（New Horizons）飞掠冥王星及其五颗卫星。

谷神星显示出活跃地质的证据，例如盐沉积与低温火山；而冥王星则拥有漂浮于氮冰冰川之上的水冰山脉，以及相当显著的大气层。谷神星显然有卤水在其地下渗流，而有证据显示冥王星可能拥有真正的地下海洋。

“黎明号”此前曾环绕小行星灶神星（Vesta）飞行。土星卫星福柏（Phoebe）曾被“卡西尼号”成像，更早之前由“旅行者2号”观测。以上三者都显示出它们曾经是矮行星的证据，对它们的探测有助于澄清矮行星的演化过程。

“新视野号”拍摄到了海王星卫星海卫一（Triton）、夸欧尔（Quaoar）、妊神星（Haumea）、阋神星（Eris）、鸟神星（Makemake），以及更小的候选体伊克西翁（Ixion）、马尼（Máni）和2014 OE394 的远距离图像[89]。夸欧尔已被提议作为中国国家航天局“申梭”双探测器的潜在飞掠目标[90]。
\subsection{相似天体}  
有一些天体在物理上与矮行星类似。这些包括：曾经的矮行星（可能仍保持流体静力学平衡的形状或存在活跃地质的证据）、行星质量卫星（满足物理定义但不满足轨道定义）以及冥王星–卡戎系统中的卡戎，在某种意义上可被视为“矮行星双星系统”。这些类别之间可能重叠，例如海卫一既是曾经的矮行星，又是行星质量卫星。
\subsubsection{曾经的矮行星}
\begin{figure}[ht]
\centering
\includegraphics[width=8cm]{./figures/6be751d540a432f2.png}
\caption{由旅行者2号拍摄图像拼接而成的海卫一（Triton）灰度镶嵌图。海卫一被认为是一颗被俘获的矮行星。} \label{fig_Dwarf_11}
\end{figure}
灶神星（Vesta）是小行星带中仅次于谷神星（Ceres）的第二大天体，曾经处于静水平衡状态，并大体呈球状；其偏离球形主要源于在其固化之后形成的两处巨型撞击盆地——Rheasilvia 与 Veneneia——所造成的巨大破坏。[91]其当前的形状参数与静水平衡状态并不相符。[92][93]

海卫一（Triton）质量大于厄里斯（Eris）与冥王星（Pluto），具有静水平衡形状，被认为是一个被海王星俘获的矮行星（很可能原本属于一个双天体系统），但它已不再直接环绕太阳运行。[94]

土卫九（Phoebe）是一个被俘获的半人马小行星，与灶神星类似，现已不再处于静水平衡，但由于早期的放射性加热作用，被认为在其历史初期曾达到过静水平衡状态。[95]
\subsubsection{行星质量卫星}
至少有十九颗卫星曾因自身重力而坍塌为完全固体（或接近固体）天体，或在历史上某个阶段通过重力自洽而达到静水平衡——尽管其中部分卫星后来已完全冻结，不再处于平衡状态。其中七颗的质量超过厄里斯或冥王星。这些大型卫星在物理性质上与矮行星并无区别，但由于它们并非直接绕太阳运行，因此不符合国际天文学联合会（IAU）的行星定义。（事实上，海卫一 Triton 是一颗被俘获的矮行星，而谷神星被认为形成于与木星和土星卫星相同的区域。）Alan Stern 将这类行星质量卫星称为“卫星行星（satellite planets）”，与矮行星和经典行星一起构成三类行星。[26]术语“planemo”（planetary-mass object）同样涵盖这三大族群。[96]

\textbf{卡戎}

关于冥王星—卡戎系统是否应被视为“双矮行星”曾存在争议。在 IAU 行星定义的草案中，冥王星与卡戎被视为一个双行星系统中的两颗行星。[19][c]IAU 目前认为卡戎不是矮行星，而是冥王星的卫星，但是否将卡戎视为矮行星这一问题可能在未来被讨论。[97]然而，卡戎是否仍保持静水平衡已不再确定。此外，轨道质心的位置不仅受两天体相对质量影响，也依赖它们之间的距离。例如：太阳—木星体系的质心位于太阳之外，但这并不意味着它们构成“双星行星系统”。因此，在正式将冥王星–卡戎定义为双矮行星之前，必须确立“双（矮）行星系统”的严格标准。
\subsection{参见}
\begin{itemize}
\item Centaur（半人马天体）\\
\item Lists of astronomical objects（天文对象列表）\\
\item List of former planets（前行星列表）\\
\item List of gravitationally rounded objects of the Solar System（太阳系内因自引力变圆的天体列表）\\
\item List of planetary bodies（行星体列表）\\
\item List of possible dwarf planets（可能的矮行星列表）\\
\item Lists of small Solar System bodies（太阳系小天体列表）
\end{itemize}
\subsection{注释}\\
a.矮行星的流体静力平衡判据无法被确认，除非有航天器对该天体进行直接探测。\\
b.以其区域内至少具有赛德娜（Sedna）质量的15个天体的最低估计值计算，见 Schwamb, Brown \& Rabinowitz (2009)[45]。\\
c.原文脚注内容如下：\\
“For two or more objects comprising a multiple object system.~…~A secondary object satisfying these conditions i.e. that of mass, shape is also designated a planet if the system barycentre resides outside the primary.~Secondary objects not satisfying these criteria are ‘satellites’.~Under this definition, Pluto’s companion Charon is a planet, making Pluto–Charon a double planet.”[19]
\subsection{参考文献}
\begin{enumerate}
\item Metzger, Philip T.; Grundy, W.\,M.; Sykes, Mark V.; 等（2022-03-01）。“卫星是行星：行星科学分类中的科学用途与文化目的论”。Icarus. 374: 114768. arXiv:2110.15285. Bibcode:2022Icar..37414768M. doi:10.1016/j.icarus.2021.114768. S2CID 240071005. 访问日期：2022-05-30。
\item “矮行星也是行星：‘新视野号’之后的行星教育学”。2021-06-27存档。
\item “深入了解 4 号灶神星（Vesta）”。NASA太阳系探索项目。2017-11-10发布。2020-02-29存档。访问日期：2020-02-29。
\item Murzi, Mauro（2007）。“科学概念的变化：什么是行星？”。科学哲学预印本（Preprint）。匹兹堡大学。2019-06-11存档。访问日期：2013-04-06。
\item Mager, Brad. “冥王星揭示（Pluto Revealed）”。discoveryofpluto.com。2011-07-22存档。访问日期：2008-01-26。
\item Ćuk, Matija; Masters, Karen（2007-09-14）。“冥王星是行星吗？”。康奈尔大学天文学系。2007-10-12存档。访问日期：2008-01-26。
\item Buie, Marc W.; Grundy, William M.; Young, Eliot F.; Young, Leslie A.; Stern, S.\,Alan（2006）。“冥王星卫星的轨道与测光：卡戎、S/2005 P1 与 S/2005 P2”。The Astronomical Journal. 132 (1): 290–298. arXiv:astro-ph/0512491. Bibcode:2006AJ....132..290B. doi:10.1086/504422. S2CID 119386667。
\item Jewitt, David; Delsanti, Audrey（2006）。《行星之外的太阳系》，载于《Solar System Update: Topical and Timely Reviews in Solar System Sciences》。Springer出版社。doi:10.1007/3-540-37683-6。ISBN 978-3-540-37683-5。2006-05-25存档。访问日期：2008-02-10。
\item Weintraub, David A.（2006）。《冥王星是行星吗？穿越太阳系的历史旅程》（Is Pluto a Planet? A Historical Journey through the Solar System）。普林斯顿大学出版社。ISBN 978-0-691-12348-6。
\item Phillips, Tony; Phillips, Amelia（2006-09-04）。“喧嚣的冥王星（Much Ado about Pluto）”。PlutoPetition.com。2008-01-25存档。访问日期：2008-01-26。
\item Brown, Michael E.（2004）。“行星的定义是什么？”。加州理工学院地质科学系。2011-07-19存档。访问日期：2008-01-26。
\item Eicher, David J.（2007-07-21）。“冥王星应被视为行星吗？”。Astronomy。2022-11-28存档。访问日期：2022-11-28。
\item “哈勃观测矮行星塞德娜（Sedna），谜团加深”。NASA哈勃望远镜网站。2004-04-14发布。2021-01-13存档。访问日期：2008-01-26。
\item Brown, Mike（2006-08-16）。“世界之战（War of the Worlds）”。纽约时报。2017-02-13存档。访问日期：2008-02-20。
\item “加州理工学院，2015-04-12检索”。2012-05-17存档。访问日期：2015-04-12。
\item “天文学家测量最大矮行星质量”。NASA哈勃望远镜网站。2007-06-14发布。2011-08-07存档。访问日期：2008-01-26。
\item Brown, Michael E. “是什么使一个天体成为行星？”。加州理工学院地质科学系。2012-05-16存档。访问日期：2008-01-26。
\item Britt, Robert Roy（2006-08-19）。“有关降级冥王星计划的细节浮现”。Space.com。2011-06-28存档。访问日期：2006-08-18。
\item “IAU关于‘行星’与‘冥族天体（plutons）’的草案定义”。国际天文学联合会。2006-08-16发布。2014-04-29存档。访问日期：2008-05-17。
\item Rincon, Paul（2006-08-25）。“投票‘被劫持’，反对声四起”。BBC新闻。2011-07-23存档。访问日期：2008-01-26。
\item Chang, Alicia（2006-08-25）。“网络商家在冥王星新闻中看到商机”。USA Today. 美联社。2008-05-11存档。访问日期：2008-01-25。
\item Brown, Michael E. “八大行星（The Eight Planets）”。加州理工学院地质科学系。2011-07-19存档。访问日期：2008-01-26。
\item “备受争议的太阳系天体获得名称”（新闻稿）。NASA。2006-09-14。2011-06-29存档。访问日期：2008-01-26。
\item Stern, Alan（2006-09-06）。“毫无畏惧地迈向第九行星（Unabashedly Onward to the Ninth Planet）”。New Horizons 网站。2013-12-07存档。访问日期：2008-01-26。
\item Wall, Mike（2011-08-24）。“冥王星行星地位的捍卫者：行星科学家Alan Stern问答”。Space.com。2012-08-14存档。访问日期：2012-12-03。
\item “大型卫星是否应被称为‘卫星行星（Satellite Planets）’？”。News.discovery.com。2010-05-14。2012-05-05存档。访问日期：2011-11-04。
\item Stern, S.\,A.; Levison, H.\,F.（2000-08-07至18）。《关于行星判据与行星分类方案的建议》（PDF）。IAU第24届大会——《Highlights of Astronomy》卷12。英国曼彻斯特（2002出版）。页205–213。Bibcode:2002HiA....12..205S。doi:10.1017/S1539299600013289。2015-09-23（PDF）存档。访问日期：2008-01-26。
\item “dwarf planet（矮行星）”。牛津英语词典（在线版）。牛津大学出版社。doi:10.1093/OED/3317852349。（需订阅或机构访问。）
\item Service, Tom（2015-07-15）。“太阳系之声：探测冥王星的预期乐章”。The Guardian。2019-12-26存档。访问日期：2019-12-26。
\item Karttunen等，编（2007）。《基础天文学》（第5版）。Springer。
\item Brown, Mike（2010）。《我如何“杀死”冥王星以及它为何咎由自取》（How I Killed Pluto and Why It Had It Coming）。Spiegel \& Grau。第223页。
\item Bailey, Mark E. “关于第5号决议的评论与讨论：行星的定义——行星的泛滥（Planets Galore）”。《Dissertatio cum Nuncio Sidereo，第三系列》——IAU2006大会官方报刊。布拉格天文研究所。2011-07-20存档。访问日期：2008-02-09。
\item “乌拉圭的Julio Fernández与Gonzalo Tancredi写入天文学史：将行星数量从9个减少为8个……以及Tancredi的注释”（西班牙语）。乌拉圭Mercedes科学与研究所。2007-12-20存档。访问日期：2008-02-11。
\item IAU（2006-08-24）。“太阳系行星的定义：决议5与6”（PDF）。IAU2006大会。国际天文学联合会。2009-06-20（PDF）存档。访问日期：2008-01-26。
\item Bowell, Edward L.\,G.; Meech, Karen J.; Williams, Iwan P.〔法语〕等（2008-12-01）。“III部：行星系统科学”。《国际天文学联合会会议录》。4 (T27A)。剑桥大学出版社：149–153。doi:10.1017/S1743921308025398。
\item “国际天文学联合会2006大会：IAU各项决议票决结果”。IAU。2006-08-24。2014-04-29存档。访问日期：2021-08-10。
\item “Plutoid 被选为类似冥王星的太阳系天体名称”。IAU。巴黎。2008-06-11。2020-11-23存档。访问日期：2021-08-10。
\item Carson, Mary Kay（2011）。《冰冷矮行星远方指南》（Far-Out Guide to the Icy Dwarf Planets）。Enslow Publishers。ISBN 9780766031876。OCLC 441945398（互联网档案馆获取）。
\item Lew, Kristi（2010）。《太空！矮行星冥王星》（Space! The Dwarf Planet Pluto）。Marshall Cavendish Benchmark。第10页。ISBN 9780761445531。OCLC 562529871（互联网档案馆获取）。
\item Darling, David（编）。“冰矮行星（Ice dwarf）”。《天体生物学、天文学与空间飞行百科全书》。2008-07-06存档。访问日期：2008-06-22。
\item “冰火山与更多发现：矮行星谷神星持续带来惊喜”。Space.com。2016年9月。2019年10月12日存档。2019年12月19日访问。
\item Castillo-Rogez, J.\,C.; Raymond, C.\,A.; Russell, C.\,T.; 等（2017-09-12）。《黎明号在谷神星：我们学到了什么？》（PDF）。行星科学与天体生物学委员会。2018年10月8日（PDF）存档。2019年10月12日访问。
\item Carroll, Michael（2019-10-23）。《谷神星：首个已知的冰矮行星》。载于Ice Worlds of the Solar System: Their Tortured Landscapes and Biological Potential。Springer Cham。doi:10.1007/978-3-030-28120-5。ISBN 978-3-030-28120-5。
\item Soter, S.（2006-08-16）。《什么是行星？》。The Astronomical Journal，132（6）：2513–2519。arXiv:astro-ph/0608359。Bibcode:2006AJ....132.2513S。doi:10.1086/508861。S2CID 14676169。
\item Schwamb, Megan E.; Brown, Michael E.; Rabinowitz, David L.（2009）。《在类赛德娜区域搜索遥远的太阳系天体》。The Astrophysical Journal，694（1）：L45–L48。arXiv:0901.4173。Bibcode:2009ApJ...694L..45S。doi:10.1088/0004-637X/694/1/L45。S2CID 15072103。
\item Margot, Jean-Luc（2015-10-15）。《定义行星的定量判据》。The Astronomical Journal，150（6）：185。arXiv:1507.06300。Bibcode:2015AJ....150..185M。doi:10.1088/0004-6256/150/6/185。S2CID 51684830。
\item Lakdawalla, Emily; 等（2020-04-21）。《什么是行星？》。planetary.org。行星学会。2022-01-22存档。2021-08-19访问。
\item Brown, Mike. “The eight planets”。gps.caltech.edu。Caltech。2011-07-19存档。访问日期：2008-01-26。
\item Jewitt, David. “Classification of Pluto”。ess.ucla.edu。UCLA。2021-08-19存档。访问日期：2021-08-19。
\item Lineweaver, Charles H.; Norman, Marc（2010）。《土豆半径：矮行星的最低尺寸》（PDF）。载于Short, W.; Cairns, I.（编）。Proceedings of 2009 Australian Space Science Conference。澳大利亚国家太空协会。页67–78。arXiv:1004.1091。ISBN 9780977574032。2023-03-10（PDF）存档。2023-08-11访问。
\item Julia Sweeney（采访者与主持）、M.\,E. Brown（受访天文学家）（2007-06-28）。《Julia Sweeney and Michael E. Brown》（播客）。Hammer Conversations。KCET。2008-06-26存档。2008-06-28访问。演员与喜剧人Julia Sweeney（《God Said Ha!》）与加州理工学院天文学家Michael E. Brown讨论使冥王星“矮化”的发现。
\item “‘行星定义’问答表”（新闻稿）。国际天文学联合会。2006-08-24。2021-05-07存档。访问日期：2021-10-16。
\item Grundy, W.\,M.; Noll, K.\,S.; Buie, M.\,W.; Benecchi, S.\,D.; Ragozzine, D.; Roe, H.\,G.（2019）。《跨海王星双星Gǃkúnǁʼhòmdímà（(229762)2007UK126）的互相轨道、质量与密度》。Icarus，334：30–38。Bibcode:2019Icar..334...30G。doi:10.1016/j.icarus.2018.12.037。S2CID 126574999。2019年4月7日（PDF）存档。
\item Nimmo, Francis; 等（2017）。《由新视野号图像确定的冥王星与卡戎的平均半径与形状》。Icarus，287：12–29。arXiv:1603.00821。Bibcode:2017Icar..287...12N。doi:10.1016/j.icarus.2016.06.027。S2CID 44935431。
\item Raymond, C.; Castillo-Rogez, J.\,C.; Park, R.\,S.; Ermakov, A.; 等（2018-09）。《黎明号数据揭示谷神星的复杂地壳演化》（PDF）。European Planetary Science Congress，第12卷。2020-01-30（PDF）存档。2020-07-19访问。
\item “根据上述定义，冥王星是‘矮行星’，并被视为新一类跨海王星天体的原型。”
\item Bruton, Dan. “绝对星等转换为小行星直径”。Stephen F. Austin州立大学物理与天文学系。2010年3月23日存档。2008年6月13日访问。
\item Ortiz, J.\,L.; Santos-Sanz, P.; Sicardy, B.; Benedetti-Rossi, G.; Bérard, D.; Morales, N.; 等（2017）。《通过恒星掩星确定矮行星Haumea的大小、形状、密度及其光环》（PDF）。Nature，550（7675）：219–223。arXiv:2006.03113。Bibcode:2017Natur.550..219O。doi:10.1038/nature24051。hdl:10045/70230。PMID 29022593。S2CID 205260767。2020-11-07（PDF）存档。2022-01-14访问。
\item Dunham, E.\,T.; Desch, S.\,J.; Probst, L.（2019-04）。《Haumea的形状、组成与内部结构》。The Astrophysical Journal，877（1）：11。arXiv:1904.00522。Bibcode:2019ApJ...877...41D。doi:10.3847/1538-4357/ab13b3。S2CID 90262114。
\item “矮行星及其系统”。行星系统命名工作组（WGPSN）。2008-07-11。2007-07-14存档。2019-09-12访问。
\item Pinilla-Alonso, Noemi; Stansberry, John A.; Holler, Bryan J.（2019-11-22）。《大型TNO表面性质：利用詹姆斯·韦伯太空望远镜扩展至更长波段的研究》。载于Dina\,Prialnik; Maria\,Antonietta\,Barucci; Leslie\,Young（编）。The Transneptunian Solar System。Elsevier。arXiv:1905.12320。
\item Dyches, Preston（2016-05-11）。《2007OR10：太阳系中最大的未命名世界》。NASA喷气推进实验室。2020-11-23存档。2019-09-12访问。
\item Porter, Simon（2018-03-27）。#TNO2018。Twitter。2018-10-02存档。2018-03-27访问。
\item 《F部“行星系统与天体生物学”年度报告2022–23》（PDF）。国际天文学联合会。2022–2023年。2023-12-08（PDF）存档。2023-12-08访问。
\item Sheppard, Scott S.; Fernandez, Yanga R.; Moullet, Arielle（2018-11-16）。《矮行星的反照率、大小、颜色及卫星：与新测量的矮行星2013FY27的比较》。The Astronomical Journal，156（6）：270。arXiv:1809.02184。Bibcode:2018AJ....156..270S。doi:10.3847/1538-3881/aae92a。S2CID 119522310。
\item Brown, Michael E. “矮行星”。加州理工学院地质科学系。2011-07-19存档。2008-01-26访问。
\item Brown, Mike. “外太阳系中有多少矮行星？”。CalTech。2011-10-18存档。2013-11-15访问。
\item Stern, Alan（2012-08-24）。《首席研究员的视角》。2014-11-13存档。2012-08-24访问。
\item Tancredi, G.; Favre, S.\,A.（2008）。《太阳系中的哪些是矮行星？》。Icarus，195（2）：851–862。Bibcode:2008Icar..195..851T。doi:10.1016/j.icarus.2007.12.020。
\item Brown, Michael（2011-08-23）。《释放矮行星！》。Mike Brown's Planets。2011-10-05存档。2011-08-24访问。
\item 在直径小于900\,km的天体中，被认为具有远高于此反照率的只有Haumea碰撞族的碎片，以及可能包括2005QU182（反照率介于0.2至0.5之间）。
\item Grundy, W.\,M.; Noll, K.\,S.; Roe, H.\,G.; Buie, M.\,W.; Porter, S.\,B.; Parker, A.\,H.; Nesvorný, D.; Benecchi, S.\,D.; Stephens, D.\,C.; Trujillo, C.\,A.（2019）。《跨海王星双星的互相轨道取向》（PDF）。Icarus，334：62–78。Bibcode:2019Icar..334...62G。doi:10.1016/j.icarus.2019.03.035。ISSN 0019-1035。S2CID 133585837。2020年1月15日（PDF）存档。2019年10月26日访问。
\item Souami, D.; Braga-Ribas, F.; Sicardy, B.; Morgado, B.; Ortiz, J.\,L.; Desmars, J.; 等（2020年8月）。《大型跨海王星天体（174567）Varda的多弦恒星掩星》。Astronomy \& Astrophysics，643：A125。arXiv:2008.04818。Bibcode:2020A\&A...643A.125S。doi:10.1051/0004-6361/202038526。S2CID 221095753。
\item Emery, J.\,P.; Wong, I.; Brunetto, R.; Cook, J.\,C.; Pinilla-Alonso, N.; Stansberry, J.\,A.; Holler, B.\,J.; Grundy, W.\,M.; Protopapa, S.; Souza-Feliciano, A.\,C.; Fernández-Valenzuela, E.; Lunine, J.\,I.; Hines, D.\,C.（2024）。《三颗矮行星的故事：来自JWST光谱的Sedna、Gonggong与Quaoar的冰与有机物》。Icarus，414：116017。arXiv:2309.15230。Bibcode:2024Icar..41416017E。doi:10.1016/j.icarus.2024.116017。
\item Rommel, F.\,L.; Braga-Ribas, F.; Ortiz, J.\,L.; Sicardy, B.; Santos-Sanz, P.; Desmars, J.; 等（2023年10月）。《通过恒星掩星测量的跨海王星天体（307261）2002MS4表面的巨大地形结构》。Astronomy \& Astrophysics，678：25。arXiv:2308.08062。Bibcode:2023A\&A...678A.167R。doi:10.1051/0004-6361/202346892。S2CID 260926329。A167。
\item Kiss, C.; Müller, T.\,G.; Marton, G.; Szakáts, R.; Pál, A.; Molnár, L.; 等（2024年3月）。《大型柯伊伯带天体（50000）Quaoar的可见光与热光曲线》。Astronomy \& Astrophysics，684：A50。arXiv:2401.12679。Bibcode:2024A\&A...684A..50K。doi:10.1051/0004-6361/202348054。
\item Cowen, R.（2007）。《特立独行的Iapetus》。Science News，第172卷，104–106页。2007年10月13日存档。
\item Thomas, P.\,C.（2010年7月）。《卡西尼提名任务后土星卫星的大小、形状及推导性质》（PDF）。Icarus，208（1）：395–401。Bibcode:2010Icar..208..395T。doi:10.1016/j.icarus.2010.01.025。2018年12月23日（PDF）存档。2015年9月25日访问。
\item Chen, Jingjing; Kipping, David（2016）。《其他世界质量与半径的概率预测》。The Astrophysical Journal，834（1）：17。arXiv:1603.08614。Bibcode:2017ApJ...834...17C。doi:10.3847/1538-4357/834/1/17。S2CID 119114880。
\item Bode, J.\,E.（编）（1801）。《1804年柏林天文年鉴》。第97–98页。2023年12月14日存档。2022年10月19日访问。
\item Slipher, V.\,M.（1930）。《跨海王星行星》。Popular Astronomy，第38卷，第415页。2021年12月11日存档。2022年10月19日访问。
\item Anderson, Deborah（2022-05-04）。《超越此世：Unicode标准批准新的天文学符号》。Unicode.org。2022年8月6日存档。2022年8月6日访问。
\item Miller, Kirk（2021-10-26）。《Unicode矮行星符号请求》（PDF）。Unicode.org。2022年3月23日（PDF）存档。2022年10月19日访问。
\item 《炼金术符号》（PDF）。Unicode.org。2022。2020年4月2日（PDF）存档。2022年10月19日访问。
\item “什么是矮行星？”。NASA喷气推进实验室。2015-04-22。2021年12月8日存档。2021年9月24日访问。
\item Miller, Kirk（2024-10-18）。《星座符号的初步展示》（PDF）。Unicode.org。2024年10月22日访问。
\item Bala, Gavin Jared; Miller, Kirk（2025-03-07）。《Phobos与Deimos符号》（PDF）。Unicode.org。2025年3月14日访问。
\item Landau, Elizabeth; Brown, Dwayne（2015-03-06）。《NASA飞船成为首个进入矮行星轨道的航天器》。NASA。2015年3月7日存档。2015年3月6日访问。
\item Verbiscer, Anne J.; Helfenstein, Paul; Porter, Simon B.; Benecchi, Susan D.; Kavelaars, J.\,J.; Lauer, Tod R.; 等（2022年4月）。《由新视野号观测的矮行星与大型KBO相位曲线的多样形状》。The Planetary Science Journal，3（4）：31。Bibcode:2022PSJ.....3...95V。doi:10.3847/PSJ/ac63a6。95。
\item Jones, Andrew（2021-04-16）。《中国将发射两艘探测器前往太阳系边缘》。SpaceNews。2021年9月29日存档。2021年4月29日访问。
\item Thomas, Peter C.; Binzelb, Richard P.; Gaffeyc, Michael J.; Zellnerd, Benjamin H.; Storrse, Alex D.; Wells, Eddie（1997）。《灶神星：来自HST图像的自转轴、大小与形状》。Icarus，128（1）：88–94。Bibcode:1997Icar..128...88T。doi:10.1006/icar.1997.5736。
\item Asmar, S.\,W.; Konopliv, A.\,S.; Park, R.\,S.; Bills, B.\,G.; Gaskell, R.; Raymond, C.\,A.; Russell, C.\,T.; Smith, D.\,E.; Toplis, M.\,J.; Zuber, M.\,T.（2012）。《灶神星的重力场及其内部结构含义》（PDF）。43rd Lunar and Planetary Science Conference（1659）：2600。Bibcode:2012LPI....43.2600A。2013年10月20日（PDF）存档。2015年7月15日访问。
\item Russel, C.\,T.; 等（2012）。《灶神星上的黎明号：原行星范式的检验》（PDF）。Science，336（6082）：684–686。Bibcode:2012Sci...336..684R。doi:10.1126/science.1219381。PMID 22582253。S2CID 206540168。2015年7月15日（PDF）存档。
\item Agnor, C.\,B.; Hamilton, D.\,P.（2006）。《海王星通过行星–双星引力相遇俘获其卫星Triton》（PDF）。Nature，441（7090）：192–194。Bibcode:2006Natur.441..192A。doi:10.1038/nature04792。PMID 16688170。S2CID 4420518。2016年10月14日（PDF）存档。2015年8月29日访问。
\item Cook, Jia-Rui C.; Brown, Dwayne（2012-04-26）。《卡西尼发现土卫Phoebe具有类行星性质》。NASA喷气推进实验室。2015年7月13日存档。
\item Basri, Gibor; Brown, Michael E.（2006）。《从星子到褐矮星：什么是行星？》（PDF）。Annual Review of Earth and Planetary Sciences，34：193–216。arXiv:astro-ph/0608417。Bibcode:2006AREPS..34..193B。doi:10.1146/annurev.earth.34.031405.125058。S2CID 119338327。2013年7月31日（PDF）存档。
\item “冥王星与太阳系”。IAU.org。国际天文学联合会。2020年4月17日存档。2013年7月10日访问。
\end{enumerate}
\subsection{外部链接}
\begin{itemize}
\item 《太阳系矮行星的可视化简介》（AnshoolDeshmukh，VisualCapitalist，2021年10月8日，图形设计：MarkBelan）
\item NPR：《矮行星或许终于能获得应有的尊重》（DavidKestenbaum，MorningEdition）
\item BBC新闻：《问答：新行星提案》，2006年8月16日
\item 《渥太华公民报》：〈反对冥王星的理由〉（P.~Surdas~Mohit），2006年8月24日
\item JamesL.~Hilton：《小行星何时成为“次要行星”？》
\item NASA：IYA2009矮行星项目
\end{itemize}